\chapter{Vertex Streams and Capabilities}
\label{ch:vertex-streams-capabilities}

This chapter is the handoff from renderer-independent graphics objects to the per-renderer Part.
It follows a draw from public vertex values through declarations and bindings, then explains what a
capability answer does---and does not---guarantee.  The distinction is essential in CNA because the
shared renderer interface contains useful defaults, explicit refusals, silent no-ops, and one
effect-discarding draw fallback.

\begin{figure}[tbp]
\centering
\begin{figurealt}{A capability query produces a reported Boolean, often inherited from a permissive default. Independent checks must then establish a usable factory or method, successful binding or draw, renderer engagement, and finally an oracle or pixel result. A failure at any stage narrows the claim.}
\resizebox{0.94\textwidth}{!}{%
\begin{tikzpicture}
  \node[cna public, text width=34mm] (query) at (-5.6,1.5) {\texttt{SupportsCapability}};
  \node[cna warning, text width=42mm] (bit) at (0,1.5) {reported Boolean\\may inherit a permissive default};
  \node[cna internal, text width=44mm] (factory) at (5.8,1.5) {factory or method exists\\and returns a usable object};
  \node[cna internal, text width=50mm] (operation) at (5.8,-1.0) {requested bind, draw, or query\\completes with the stated semantics};
  \node[cna evidence, text width=31mm] (engaged) at (0,-1.0) {renderer engaged};
  \node[cna evidence, text width=31mm] (oracle) at (-5.6,-1.0) {oracle or pixels\\compared};
  \draw[cna flow] (query) -- (bit);
  \draw[cna flow] (bit) -- (factory);
  \draw[cna flow] (factory) -- (operation);
  \draw[cna flow] (operation) -- (engaged);
  \draw[cna flow] (engaged) -- (oracle);
\end{tikzpicture}}
\end{figurealt}
\caption{Capability reporting and capability evidence. A true query result is an API response,
not a substitute for a usable factory, successful operation, or discriminating oracle.}
\label{fig:capability-evidence}
\end{figure}


\section{Public vertex objects are not upload layouts}

The built-in public types preserve XNA-shaped fields and \cnaclass{IVertexType} behavior, but their
C\texttt{++} object representation is not necessarily the GPU byte representation.  In particular,
\cnaclass{Color} and the vertex interfaces are polymorphic.  Chapter~\ref{ch:geometry-color-layout}
documents the plain internal stream structs and their asserted 16/20/24/32/48/52/68-byte layouts.

Typed \cnaclass{VertexBuffer::SetData} and \cnaclass{GraphicsDevice::DrawUserPrimitives} routes
therefore pack supported public values into those stream forms.  Raw-pointer routes are a different
contract: the caller supplies already-packed bytes and, where required, an explicit
\cnaclass{VertexDeclaration}.  Passing an array of public objects through \texttt{void*} does not
turn their vptr-bearing C\texttt{++} layout into the XNA stride.

\section{VertexDeclaration is the byte contract}

A \cnaclass{VertexDeclaration}'s layout is immutable after construction. Each element stores an offset,
usage, and usage index.  Its \cnaclass{VertexElementFormat} describes the byte shape; the
declaration's stride is either explicit or
computed as the largest \(\text{offset}+\text{format size}\). The public default constructor is a
deliberate exception: it creates an empty, zero-stride declaration used by the legacy
\cnaclass{VertexBuffer}(device, count) convenience path. The constructors that explicitly receive
an element list reject an empty list with \cnaclass{ArgumentNullException}; explicit non-positive
strides throw the out-of-range exception.

The declaration reaches every renderer through a mandatory call.  The vertex-buffer renderer
interface method is \texttt{SetVertexDeclaration}.%
\index{IVertexBufferRenderer!SetVertexDeclaration} An old description that says CNA
selects layouts from stride alone is obsolete.  Stride is still visible because several renderer
implementations optimize known built-in layouts, but a faithful path must also validate and bind
the declared elements.

\cnaclass{VertexElement} itself is mutable and its hash currently returns zero, following the
corresponding FNA limitation.  Do not use the hash as a declaration identity; compare the complete
element sequence and stride when building a cache key.

\section{Buffers, index widths, and dynamic updates}

\cnaclass{VertexBuffer} stores a declaration, vertex count, usage, and renderer-owned resource.
Its typed transfer families cover CNA's seven built-in stream shapes; the CNAEXT raw route accepts
a pointer, element count, and stride.  Copy is deleted and move is explicit/noexcept, which avoids
duplicating a native buffer handle.

\cnaclass{IndexBuffer} exposes 16-bit and 32-bit transfer forms selected by
\cnaclass{IndexElementSize}.  The semantic choice is not the enum's ordinal.  Content and oracle
tools must compare the 16/32 meaning rather than assuming an underlying 0/1 value, a mismatch that
has produced real cross-implementation disagreements.

\cnaclass{DynamicVertexBuffer} and \cnaclass{DynamicIndexBuffer} add a content-lost surface and
\texttt{SetData(..., SetDataOptions)}.  At the pin, the dynamic update destination always begins at
offset zero; the options influence native update strategy, not an arbitrary destination offset.
In normal CNA operation, their content-lost property remains false.

\section{Bindings carry offset and input rate}

\cnaclass{VertexBufferBinding} combines a buffer, \texttt{VertexOffset}, and
\texttt{InstanceFrequency}.  Frequency zero is per-vertex input; a positive frequency is
per-instance input.  \cnaclass{GraphicsDevice} accepts up to sixteen binding slots through
\texttt{SetVertexBuffers}.\index{GraphicsDevice!SetVertexBuffers} It permits null buffers as
unused slots.  The public ceiling is not proof that the selected
renderer can consume sixteen active streams.

The classic shapes need no multi-stream capability:

\begin{itemize}
  \item one per-vertex binding;
  \item one per-vertex binding plus one per-instance binding.
\end{itemize}

More than one per-vertex stream, more than one per-instance stream, or a binding count beyond the
renderer-reported native ceiling invokes the \texttt{MultiStreamVertexInput} gate.  A false answer
causes \cnaclass{NotSupportedException} before native submission instead of letting the renderer
silently read only stream zero.

Streams compose by \((\text{VertexElementUsage},\text{UsageIndex})\), not merely by slot order. The
shared layer claims each semantic pair while walking bindings. A later stream whose declaration is
entirely duplicated contributes nothing, matching the FNA3D driver rule that such a stream is not in
use. A declaration that repeats only some earlier semantics is rejected: CNA's combined-stride
description cannot express a record whose fields are partly suppressed. Split declarations should
therefore use distinct semantic/index pairs deliberately.

\subsection{The minimum-offset fold}

For multiple per-vertex bindings, CNA finds the smallest \texttt{VertexOffset} and folds it into the
draw's \texttt{vertexStart} or \texttt{baseVertex}.  Each binding carries only its non-negative
remainder.  Renderers already multiply the shared element-unit base by each stream's own stride,
so this scheme applies the common portion exactly once without requiring equal strides.

With one stream, the transformation is byte-identical to the older behavior.  With several streams,
folding each full offset independently into one global base would be wrong; dropping the remainder
would align only the earliest stream.  This normalization belongs to the shared layer so every
renderer sees the same binding geometry.

\section{Draw validation has a deliberate order}

Every buffer-backed 3D draw first asks the renderer to enforce its 3D policy, then requires the
necessary vertex/index buffers and a currently applied \cnaclass{Effect}.  Missing effect state is a
plain \texttt{std::runtime\_error}, not an XNA-shaped exception.

Range validation occurs before renderer capability validation.  An index span leaving the index
buffer or a declared vertex range leaving the vertex buffer is invalid on every renderer and must
produce the same public argument failure even on a 2D-only product.  Only after those universal
invariants are established does CNA reject an unsupported stream shape or instancing path.

The user-draw overloads follow the same conceptual stages: establish a declaration, pack typed
objects when necessary, validate counts and ranges, construct transient upload data, fill
\cnaclass{GpuDrawParams} from the applied effect, and enter the renderer contract.

\section{Thirteen capability questions}

\cnaclass{GraphicsCapability} contains thirteen entries at the pinned revision:

\begin{center}
\small
\begin{tabularx}{\textwidth}{lX}
\toprule
Capability & Exact question \\
\midrule
\texttt{ThreeD} & Is the vertex/index/effect 3D pipeline supported? \\
\texttt{DepthStencilBuffer} & Is a complete real depth/stencil attachment available? \\
\texttt{MultiSampleAntiAliasing} & Can a sample count above one be applied? \\
\texttt{MultipleRenderTargets} & Can more than one target be bound simultaneously? \\
\texttt{AnisotropicFiltering} & Does the current device/driver honor anisotropic filtering? \\
\texttt{WireFrame} & Is \texttt{RasterizerState.FillMode=WireFrame} real? \\
\texttt{OcclusionQuery} & Is there a real query with Begin/End/completion/pixel result? \\
\texttt{CustomEffects} & Can a non-stock effect enter the documented custom-effect path? \\
\texttt{Texture3D} & Does real volume texture storage and transfer exist (not necessarily shader sampling)? \\
\texttt{MultiStreamVertexInput} & Can repeated input rates be represented across several bindings? \\
\texttt{Instancing} & Can \texttt{DrawInstancedPrimitives} submit distinct instance records? \\
\texttt{StencilBuffer} & Is a real stencil plane available independently of a depth/3D claim? \\
\texttt{AdditiveBlending} & Does additive compositing avoid silent source-over degradation? \\
\bottomrule
\end{tabularx}
\end{center}

These are CNA contract questions, not a dump of native API feature bits.  For example, Skia can
report bounded \texttt{Texture3D} CPU storage while its general 3D pipeline remains false.  GDI can
offer a 2D stencil mask without claiming a depth attachment.  HTML DOM additive support depends on
whether the browser honors CSS \texttt{plus-lighter}.

\section{The default polarity is permissive}

\texttt{IGraphicsRenderer::SupportsCapability} returns true by default for every entry except
\texttt{MultiStreamVertexInput}; \texttt{StencilBuffer} delegates to the older stencil query.
This opt-out design protects mature renderers from boilerplate, but a new or narrow renderer that
forgets an override overclaims.

The implementation is the authority.  At the pin, HEADLESS and SOFTWARE report MRT through the
default even though their bind paths do not provide simultaneous target rendering.  WEBGPU reports
MRT and occlusion queries while rejecting a second target and returning no real query object;
SDL\_GPU similarly inherits a positive query answer without a query implementation. Custom effects
add two direct false positives: WEBGPU inherits true while retaining the null effect factory, and
BGFX reports true on a real renderer even though its returned effect object's
\texttt{CompileProgram()} always returns false. FNA3D is the useful opposite example: it executes
six closed stock bytecode programs but truthfully reports custom effects false
(Chapter~\ref{ch:shader-fx-gap}). Appendix \ref{ch:feature-matrix} records such contradictions
rather than copying the Boolean as support.

A capability check is therefore necessary but, for known contradictory entries, not sufficient.
Pair it with a resource/draw path and an observable oracle until the upstream report becomes
truthful.

\section{Unsupported 3D policy is separate from capability}

Each \cnaclass{GraphicsDevice} carries
\cnaclass{Unsupported3DGraphicsCallBehavior}, with two values:

\begin{description}
  \item[\texttt{Throw}] The default.  A deliberately unsupported 3D call follows the renderer's
        established refusal path.
  \item[\texttt{WarnAndStub}] The renderer logs once per method name and returns a safe no-op or
        null-object resource where the contract defines one.
\end{description}

Changing the policy clears the warn-once history.  It does not change capability answers and does
not suppress invalid arguments, disposed resources, driver failures, or unfinished work on a
renderer that otherwise claims 3D.  There is no environment variable for it; applications opt in
through the device API.

Six deliberately 2D-focused families override the early \texttt{Ensure3DSupported} hook:
Blend2D, Direct2D, HTML DOM, OpenVG, Skia, and SVG DOM.  Other 2D/refusal behavior can live in
individual factory or draw methods, which is why an absent override is not itself proof of 3D.

\section{Four interface-default failure shapes}

The shared renderer contract is over two thousand lines and not every default has the same
semantics:

\begin{enumerate}
  \item \textbf{Effect-discarding fallback.}\par
        The ordinary effect-aware draw defaults discard the \cnaclass{GpuDrawParams} block, then
        call coloured primitive submission.  Texture, lighting, fog,
        skinning, and custom effect state can disappear without an exception.
  \item \textbf{Hard refusal.} Instancing and backbuffer readback defaults throw unless a narrowly
        defined unsupported-3D policy path intercepts the operation.
  \item \textbf{Silent no-op.} Several state/debug hooks have empty defaults; a public cache may
        update even when no native state changes unless the shared setter guards that path.
  \item \textbf{Null factory.} Optional resource/effect factories may return null, leaving the
        wrapper to refuse, degrade, or expose an inert object.
\end{enumerate}

The renderer chapters in Part~\ref{part:renderers} classify the actual destination of each family instead of inferring it
from override presence.  Thirty-one families own both ordinary extended draw methods; eleven
inherit both.  Of the inherited group, DIRECTX10 alone reaches a real coloured draw after effect
state is discarded, while refusal/no-op/trace families terminate differently.  Four additional
override families retain a conditional coloured tail.

\section{A portable draw claim}

“The triangle rendered” is incomplete.  A durable vertex-path claim names:

\begin{itemize}
  \item the public identity and implementation family;
  \item the vertex declaration, packed stride, index width, and binding rates;
  \item the capability answers and any known report/path contradiction;
  \item the applied effect and whether the renderer consumed its \cnaclass{GpuDrawParams};
  \item the native driver/host that engaged;
  \item the pixel, query, or structured oracle that distinguished the expected route from a
        coloured fallback or no-op.
\end{itemize}

With those nouns present, a failure can be localized to packing, binding normalization, capability
reporting, effect dispatch, or native rendering instead of being summarized as “this renderer does
not work.”
